% Chapter 3 -- Requirements, Impacts and Constraints
% Template-required chapter (CSE400 Fall 2024 onwards)

\section{Final Specifications and Requirements}

\lecseg{} is a text-first, offline lecture-video topic segmentation system.
Its concrete technical requirements are:

\begin{enumerate}
  \item \textbf{Input:} A Whisper transcript (JSON) or pre-chunked auto-captions
        from yt-dlp. Minimum viable input is a plain-text transcript chunked into
        pseudo-sentences of 20--30 words.
  \item \textbf{Embeddings:} One sentence-transformer model from the supported set
        (BGE-M3, E5-large-v2, MPNet, MiniLM). BGE-M3 is the default; CPU inference
        requires $\approx$3 minutes per 60-minute lecture on a mid-range machine.
  \item \textbf{Output:} A list of boundary timestamps (seconds) with optional
        LLM-generated chapter titles when Ollama is available.
  \item \textbf{Ground-truth format:} Chapter timestamps in the \texttt{data/gt/}
        JSON schema; used only at evaluation time, not at inference time.
  \item \textbf{Software stack:} Python 3.10+, PyTorch $\geq$2.0, sentence-transformers,
        spaCy \texttt{en\_core\_web\_sm}, scikit-learn, librosa. All packages are
        MIT- or Apache-2.0-licensed.
  \item \textbf{Hardware:} No GPU required for inference. Whisper large-v3
        transcription benefits from a GPU but can run on CPU for short videos.
\end{enumerate}

\section{Societal Impact}

Lecture segmentation directly lowers the barrier to self-directed learning.
When a two-hour recorded lecture is automatically divided into labelled chapters,
a student who has ten minutes can locate and revisit a specific concept rather
than scrubbing through an undivided recording.
At scale this has measurable equity implications: learners without institutional
library access, returning students balancing work and study, and learners in
lower-bandwidth environments all benefit from shorter, addressable clips rather
than monolithic recordings.

The \lecseg{} pipeline is fully offline and transmits no student data to any
third-party service, which preserves learner privacy and allows deployment in
institutions with strict data-residency requirements.
Because all components are open-source, the system can be self-hosted at zero
licensing cost by any university, eliminating the access barrier that commercial
captioning and chapter-generation services create.

\section{Environmental Impact}

The pipeline is designed to run on a single CPU-only consumer laptop.
The most computationally intensive inference step, BGE-M3 sentence embedding,
processes a 60-minute lecture in under three minutes on a mid-range CPU. No GPU
is required for inference or evaluation.

Transcription was performed once on a rented cloud GPU (RTX 5090, $\approx$1.5
GPU-hours total for all 30 videos), after which transcripts are cached and reused
indefinitely.
The method selector (logistic regression over hand-crafted features) has
negligible training energy.
Contrast this with transformer fine-tuning approaches: a single training run of
a 110M-parameter BERT model for 3 epochs on comparable data typically consumes
$>$50 GPU-hours.
The \lecseg{} design therefore has a substantially lower carbon footprint than
fine-tuning-based segmentation systems.

\section{Ethical Issues}

All 30 \dataset{} videos are publicly available on YouTube under the channel
owners' terms of service.
No personally identifiable information beyond publicly available lecture content
is stored or redistributed.
The annotation labels capture only temporal boundary positions; no judgements
about individual lecturers' teaching quality or competence are expressed.

The LLM draft annotations (Ollama/llama3.1:8b) run entirely on-device; no
prompts or transcripts are transmitted to external services.
The system cannot and does not make decisions about student grades, access, or
academic progression. It is a navigation aid, not an evaluative tool.
Researchers reusing \dataset{} should note that creator-supplied YouTube chapter
labels reflect editorial intent, not authoritative pedagogical ground truth;
treating them as such would be a misuse of the corpus.

\section{Standards}

The evaluation protocol follows community-accepted standards for text
segmentation:
\begin{itemize}
  \item $P_k$~\cite{beeferman1999statistical} and WindowDiff~\cite{pevzner2002critique}
        are the de-facto standard metrics for segmentation evaluation, used in
        all major benchmarks since 2000.
  \item Bootstrap confidence intervals (1{,}000 resamples, 95\%) follow
        \cite{efron1993_bootstrap}.
  \item Holm-corrected Wilcoxon signed-rank tests follow the standard correction
        for family-wise error rate in pairwise comparisons.
  \item The FAIR principles (Findable, Accessible, Interoperable, Reusable) are
        met: the corpus, labels, code, and pre-computed results are all published
        in the public GitHub repository under permissive licences.
\end{itemize}

\section{Project Management Plan}

The project followed a staged milestone structure over 20 weeks:

\begin{center}
\begin{tabular}{lll}
  \toprule
  Phase & Weeks & Deliverable \\
  \midrule
  1. Dataset construction \& annotation & 1--6  & Annotated \dataset{} corpus \\
  2. Baseline implementation \& ablation & 7--10 & Ablation results table \\
  3. Novel method development (N1--N4)   & 11--14 & Component code + eval \\
  4. Statistical testing \& external validation & 15--17 & Results JSON + charts \\
  5. Thesis writing \& revision          & 18--20 & Final thesis PDF \\
  \bottomrule
\end{tabular}
\end{center}

Each phase produced a concrete, independently reviewable artifact.
Version control (Git) was used throughout; the full repository history at
\url{https://github.com/IsmamK/Thesis_TopicSegmentation} serves as the
project log.
Task tracking followed a numbered task list (Tasks T01--T47) maintained in
\texttt{progress.yaml}.

\section{Risk Management}

Three principal risks were identified at project start; all materialised and
were mitigated:

\begin{enumerate}
  \item \textbf{Annotation bottleneck.}
        \textit{Risk:} Manual double-annotation of 32 hours of video would take
        $>$40 annotator-hours, creating a schedule bottleneck.
        \textit{Mitigation:} LLM draft annotations (Ollama llama3.1:8b) were used
        as a first pass, reducing human effort to review and correction
        ($\approx$8 hours total).

  \item \textbf{Transcription cost.}
        \textit{Risk:} Running Whisper large-v3 locally on 32.52 hours of audio
        would require $>$48 CPU-hours.
        \textit{Mitigation:} A short GPU burst on a rented cloud instance (RTX 5090,
        1.5 hours, $\approx$\$1 USD) completed all transcription in one session.

  \item \textbf{Benchmark data leakage.}
        \textit{Risk:} Using any of the 30 \dataset{} videos in the external
        validation set would inflate reported generalisation.
        \textit{Mitigation:} A hard exclusion list (\texttt{LECSEG30\_IDS}) is
        maintained in \texttt{scripts/external\_eval.py} and checked at runtime.
\end{enumerate}

\section{Economic Analysis}

\textbf{Development cost:}
\begin{itemize}
  \item Cloud GPU rental (transcription): $\approx$\$1 USD (one-time)
  \item Software licences: \$0 (all components MIT/Apache-2.0)
  \item Storage: $\approx$2 GB for embeddings and transcripts (free on any local machine)
\end{itemize}

\textbf{Deployment cost per university:}
\begin{itemize}
  \item Transcription: $\approx$\$0.03 per lecture-hour using the Whisper API,
        or free using local CPU (slower).
  \item Inference (embedding + segmentation): $\approx$3 CPU-minutes per
        lecture-hour; effectively free on any server.
  \item Storage: $\approx$4 MB per lecture-hour for embeddings.
\end{itemize}

\textbf{Cost-benefit:}
A university with 1{,}000 recorded lectures ($\approx$2{,}000 hours) could
auto-segment the entire archive for under \$60 in transcription API costs, or
free using local CPU over approximately one week.
The navigability benefit (students finding content in minutes rather than
scrubbing through hours) represents a measurable reduction in study-time
waste across the institution.
Compared to commercial lecture-segmentation SaaS products (typical cost:
\$0.50--\$2.00 per lecture-hour), the open-source \lecseg{} pipeline offers a
cost saving of 95\%+ for institutions willing to self-host.

